\begin{abstract}
The kissing-number problem is a classical problem in discrete geometry
whose exact solution is known in only a few dimensions. Recent
artificial-intelligence approaches have begun to discover improved
configurations through large-scale numerical and combinatorial search,
but converting such searches into general mathematical constructions and
rigorous proofs remains challenging. Here we use Qiushi Engine, an autonomous
multi-agent research system, to investigate kissing numbers and obtain
new lower bounds in nineteen dimensions: $25$, $27$, $32$--$39$, $43$,
$45$, and $49$--$55$. The resulting constructions arise from distinct
structural mechanisms, including coordinated motions of contact layers,
labelled direction reuse, joint support exchanges, signed-code
replacements, cross-shell lattice constructions, low-overlap lattice
isometries, and spherical-design moment certificates. They yield sharp
capacities for parameterized signed-code models, deterministic image-union
guarantees, and exact section and projection counts controlled by embedded
root systems and anchor-graph statistics. These methods
yield, among others, $K(25)\ge197580$, $K(27)\ge201567$,
$K(38)\ge591900$, $K(43)\ge2553792$, $K(45)\ge7380090$, and
$K(55)\ge53301140$. The autonomous system carried out the construction
searches, mathematical analysis and computational verification, while
each final result was reduced to explicit mathematical arguments and
independently checkable finite certificates. Our results illustrate how
autonomous research systems can move beyond optimization within fixed
formulations to discover new mathematical representations and
constructions at scale.

\par\smallskip\noindent\textbf{Keywords}\quad
Autonomous mathematical discovery; kissing number; geometric deformation;
constant-weight code; lattice lifting; signed design; spherical design
\end{abstract}
